% §7.5 Worked Application: Bloom Filter Diagnostics
%
% Cross-reference resolution status:
%
% DEFINED here:
%   \label{sec:bloom-diagnostics}                 -- this section
%
% ALREADY DEFINED (chapter 5):
%   \Cref{sec:bloom-from-scratch}                 -- chapter 5 §5.2
%   \Cref{def:bloom-filter}                       -- chapter 5 §5.2
%   \Cref{eq:bloom-fpr}                           -- chapter 5 §5.2
%   \Cref{thm:bloom-optimal-k}                    -- chapter 5 §5.2
%
% ALREADY DEFINED (chapter 7, this plan):
%   \Cref{sec:binomial-estimator}                 -- §7.2
%   \Cref{sec:confidence-intervals}               -- §7.3
%   \Cref{def:wilson-ci}                          -- §7.3
%   \Cref{sec:derived-measures}                   -- §7.4
%   \Cref{def:ppv}                                -- §7.4 PPV definition
%   \Cref{thm:delta-method}                       -- §7.4 delta-method theorem
%   \Cref{prop:ppv-variance}                      -- §7.4 PPV variance via delta method
%
% §7.5 applies the full chapter-7 toolkit (binomial estimator from §7.2,
% Wilson CI from §7.3, PPV definition and delta-method CI from §7.4) to
% one concrete scenario: a deployed Bloom filter with measured rates on
% a workload with measured prevalence. The diagnostic question at the
% end forks: re-tune the filter, or wait for higher prevalence?

\section{Worked Application: Bloom Filter Diagnostics}
\label{sec:bloom-diagnostics}

The chapter-7 toolkit is now complete: \Cref{sec:binomial-estimator}
proved the empirical-proportion estimator is unbiased and efficient;
\Cref{sec:confidence-intervals} replaced the Wald interval with Wilson
as the default; \Cref{sec:derived-measures} carried both into PPV via
the delta method. The remaining task is to use all three on a single
deployment problem and watch the diagnostic question fall out of the
numbers. The scenario below is a spam filter implemented as a Bloom
filter, the running example of \Cref{sec:bloom-from-scratch}; the
deployment question is whether the filter is precise enough on the
workload at hand, and if not, what to change.

\subsection*{Scenario}

A spam filter is implemented as a Bloom filter
(\Cref{def:bloom-filter,sec:bloom-from-scratch}) with parameters
$(m, k, n) = (50000, 7, 5000)$. The bit-to-item ratio $m/n = 10$ and
$k = 7$ put the filter in the optimal regime of
\Cref{thm:bloom-optimal-k}. Applying \eqref{eq:bloom-fpr} with
$kn/m = 0.7$,
\[
  \fprate
  \;\approx\; \bigl(1 - e^{-0.7}\bigr)^7
  \;\approx\; (0.5034)^7
  \;\approx\; 0.0082.
\]
The construction guarantees $\fnrate = 0$: no spam message is missed.
A workload of $N = 10000$ messages arrives, a mix of spam and
not-spam in proportions the deployment does not control. The pair
$(\fprate, \fnrate) = (0.0082, 0)$ describes the filter; the
workload's prevalence is a separate quantity that does not enter the
filter at all. The PPV at deployment is a joint property of both, as
\Cref{sec:derived-measures} insisted. The deployment question is
what fraction of $\top$-flagged messages are actually spam, and what
to do if that fraction is not high enough for the operational
requirement.

\subsection*{Empirical observation}

Ground truth is available on this workload (the operator labels a
sample, or downstream feedback eventually arrives). Of the $N = 10000$
messages, $n_\top = 1000$ are spam and $n_\bot = 9000$ are not. The
filter reports $\top$ on every spam message by $\fnrate = 0$ and on
$\mathrm{FP} = 92$ of the non-spam. The empirical false-positive rate
is therefore
\[
  \hat{\fprate}
  \;=\; \frac{\mathrm{FP}}{n_\bot}
  \;=\; \frac{92}{9000}
  \;\approx\; 0.01022,
\]
about twenty-five percent above the theoretical $0.0082$. The Wilson
$95\%$ confidence interval (\Cref{def:wilson-ci}) at $z = 1.96$ has
center
\[
  \frac{0.01022 + 1.96^2 / 18000}{1 + 1.96^2 / 9000}
  \;\approx\; \frac{0.01044}{1.000427}
  \;\approx\; 0.01043
\]
and half-width
\[
  \frac{1.96}{1.000427}
    \sqrt{\frac{0.01022 \cdot 0.98978}{9000} + \frac{1.96^2}{4 \cdot 9000^2}}
  \;\approx\; 1.959
    \sqrt{1.124 \times 10^{-6} + 1.186 \times 10^{-8}}
  \;\approx\; 0.00209,
\]
giving the interval $[0.00834, 0.01252]$. The theoretical
$\fprate = 0.0082$ sits just outside the lower endpoint. The
discrepancy is real and not an artifact of small samples: at
$n_\bot = 9000$ the Wilson interval is narrow enough to detect a
$25\%$ deviation, and the theoretical value falls below the lower
endpoint of the interval by roughly one standard error. Probable
causes include real-world hash collisions that mildly violate the
uniform-hash assumption behind \eqref{eq:bloom-fpr}, or an item count
slightly above the design $n = 5000$ pushing the filter into a less
favourable regime. The diagnostic point is that the deployment
$\fprate$ is an empirical quantity to be measured, not a design
parameter to be trusted on faith. We carry $\hat{\fprate}$ forward in
the calculations below; the theoretical value $0.0082$ is the
designer's promise, and the empirical $0.0102$ is what the workload
actually sees.

\subsection*{PPV computation}

The workload's prevalence point estimate is
$\hat\pi = n_\top / N = 0.1$. Combining this with $\hat{\fprate}$ and
$\hat{\fnrate} = 0$ in \Cref{def:ppv} gives
\[
  \widehat{\PPV}
  \;=\; \frac{0.1 \cdot 1}{0.1 \cdot 1 + 0.9 \cdot 0.01022}
  \;=\; \frac{0.1}{0.10920}
  \;\approx\; 0.9158.
\]
About $92\%$ of $\top$-flagged messages are actually spam. A
delta-method confidence interval follows from
\Cref{prop:ppv-variance}. Conditioning on $n_\top, n_\bot$ and
substituting $\bar t = n_\top \cdot 1 = 1000$, $\bar f = n_\bot
\hat{\fprate} = 92$, $\sigma_t^2 = n_\top \cdot 1 \cdot 0 = 0$, and
$\sigma_f^2 = n_\bot \hat{\fprate}(1 - \hat{\fprate}) \approx 91.06$,
\[
  \Var[\widehat{\PPV}]
  \;\approx\; \frac{1000^2 \cdot 91.06}{(1000 + 92)^4}
  \;\approx\; 6.40 \times 10^{-5},
\]
so the standard error is about $0.00800$ and the Wald-style
delta-method $95\%$ interval is
\[
  \widehat{\PPV} \;\pm\; 1.96 \cdot 0.00800
  \;\approx\; 0.9158 \pm 0.0157
  \;\approx\; [0.900, 0.932].
\]
The interval is informative: PPV is bounded above by about $0.93$ at
this workload's prevalence and at this filter's measured $\fprate$,
with $95\%$ confidence on the sampling distribution. The
boundary-failure concern of \Cref{sec:confidence-intervals} does not
apply here because $\widehat{\PPV}$ is far from $0$ and $1$, and at
the conditional sample sizes $n_\top = 1000, n_\bot = 9000$ the Normal
approximation underlying the delta method is comfortable. The variance
formula is also simplified by $\fnrate = 0$: the $\sigma_t^2$ term
vanishes and the entire uncertainty in PPV comes from the
false-positive count. Were $\fnrate > 0$, the recall would itself be
a random quantity and both terms in \Cref{prop:ppv-variance} would
contribute.

\subsection*{Diagnostic question}

Suppose the operational requirement is $\PPV > 0.95$: at most one
$\top$-report in twenty should be wrong. The point estimate $0.9158$
falls short. The interval $[0.900, 0.932]$ excludes $0.95$. The
deployment needs a fix.

\Cref{def:ppv} has only three knobs: $\fprate$, $\fnrate$, and $\pi$.
The filter has $\fnrate = 0$ already. Two routes are open.

\smallskip
\noindent\emph{Route A: re-tune.} Lower $\fprate$ by enlarging the bit
array or shifting $k$. To match $\PPV = 0.95$ at $\pi = 0.1$,
\[
  0.95 \;=\; \frac{0.1}{0.1 + 0.9 \fprate}
  \quad\Longleftrightarrow\quad
  \fprate \;=\; \frac{0.1 \cdot 0.05}{0.9 \cdot 0.95}
  \;\approx\; 0.00585.
\]
At $\fprate = 0.005$ the formula gives $\PPV = 0.1 / (0.1 + 0.9 \cdot
0.005) \approx 0.9569$, comfortably above the target. Reaching
$\fprate = 0.005$ from $0.01022$ by \eqref{eq:bloom-fpr} costs about
$30\%$ more bits at fixed $k$, or a small re-balancing of $k$ at
fixed $m$. The $30\%$ space cost is a one-time hardware decision and
buys an $\fprate$ that survives the worst-case workload.

\smallskip
\noindent\emph{Route B: wait.} Keep the filter as-is and bound the
acceptable workload. At fixed $\fprate = 0.01022$, the prevalence
required for $\PPV = 0.95$ is
\[
  0.95 \;=\; \frac{\pi}{\pi + (1 - \pi) \cdot 0.01022}
  \quad\Longleftrightarrow\quad
  \pi \;\approx\; 0.163.
\]
At $\pi = 0.2$ the formula gives $\PPV = 0.2 / (0.2 + 0.8 \cdot
0.01022) \approx 0.9607$, also above the target. Doing nothing
structural is acceptable when the spam rate naturally rises above
about one in six; below that, the target is missed. This route
trades capital cost for the assumption that the workload mix will
shift in the desired direction, which the deployment does not get to
choose. It is the right route when prevalence is expected to climb
for independent reasons, and the wrong route when prevalence might
stay below the threshold indefinitely.

\smallskip
The choice is operational rather than statistical. Route~A spends
storage to lower $\fprate$; Route~B accepts a residual $8\%$ false
share until the workload's prevalence rises on its own. Both
calculations used the same three formulas: \eqref{eq:bloom-fpr} for
the filter's $\fprate$, the Wilson interval for its empirical
estimate, and \Cref{def:ppv} for the translation into precision at
workload prevalence. What the chapter-7 toolkit added beyond the
chapter-5 design formulas is the honest accounting of uncertainty. A
point estimate of PPV alone, $0.9158$, would have suggested the
target $0.95$ was nearly met and a small tweak would suffice. The
delta-method interval $[0.900, 0.932]$ shows that the data rule out
$0.95$ at $95\%$ confidence: a small tweak would not have sufficed.
The diagnostic question reduces to one division, run forward and
backward, but the conclusion that a fix is needed at all rests on the
interval, not the point.
